% !TEX root = ../../main.tex
% C3.6 — Tìm kiếm vector (RÚT GỌN 2026-09-29; bản cũ ở report/archive/)
% Nguồn: HNSW (malkov2020hnsw, TPAMI 42(4)); Milvus (wang2021milvus, SIGMOD 2021); tài liệu Milvus HNSW + filtered search.
% Khớp app: storage/milvus/vectors.py (HNSW, IP, M=16, efConstruction=128, ef=64; filter area + READY + camera_ids + khoảng thời gian);
% services/track_search.py (top_k ∈ {4,8,12,16}; đối chiếu PostgreSQL, bỏ kết quả không hợp lệ, nhân đôi limit, tối đa 3 vòng → limit ≤ 64 = ef).
\section{Tìm kiếm vector}
\label{sec:3.6}

\subsection{Tìm kiếm lân cận gần nhất và chỉ mục HNSW}
\label{sec:3.6.1}
\label{sec:3.6.2}

Khi mỗi lần xuất hiện được đại diện bởi một vector (Mục~\ref{sec:3.5}) và truy vấn cũng được mã hóa thành một vector $\mathbf{q}$, việc tìm kiếm trở thành bài toán \emph{tìm $k$ lân cận gần nhất}: trong $n$ vector đã lưu, tìm $k$ vector có tích vô hướng với $\mathbf{q}$ lớn nhất, tương đương độ tương đồng cosine khi vector đã chuẩn hóa. Tìm kiếm vét cạn luôn chính xác nhưng chi phí tăng tuyến tính theo $n$, nên thời gian phản hồi tăng dần khi dữ liệu tích lũy. \emph{Tìm kiếm lân cận gần nhất gần đúng} (Approximate Nearest Neighbor, ANN) xây dựng trước một chỉ mục để khi truy vấn chỉ cần xét một phần nhỏ dữ liệu, chấp nhận bỏ sót một phần nhỏ lân cận thực sự để đổi lấy tốc độ.

HNSW (Hierarchical Navigable Small World)~\cite{malkov2020hnsw} là phương pháp ANN dựa trên đồ thị: mỗi vector là một đỉnh nối với một số vector gần nó, và việc tìm kiếm lặp lại việc di chuyển sang đỉnh kề gần truy vấn hơn. Các đỉnh được tổ chức thành nhiều tầng, tầng trên thưa để di chuyển nhanh tới vùng gần truy vấn, tầng dưới cùng chứa mọi vector để tìm chính xác trong vùng đó, nên độ phức tạp tìm kiếm tăng theo hàm logarit của số phần tử~\cite{malkov2020hnsw}. Chỉ mục có ba tham số chính~\cite{milvus2026hnsw}: số cạnh tối đa của mỗi đỉnh $M$ và số ứng viên khi xây dựng chỉ mục \textit{ef\-Con\-struc\-tion}, càng lớn thì đồ thị càng tốt nhưng tốn bộ nhớ và thời gian xây dựng hơn; và số ứng viên khi tìm kiếm $\mathit{ef}$, càng lớn thì càng chính xác nhưng chậm hơn, và không nên nhỏ hơn số kết quả cần lấy.

\subsection{Tìm kiếm vector kết hợp điều kiện lọc}
\label{sec:3.6.3}

Truy vấn trong hệ thống giám sát thường kèm điều kiện, như chỉ trong khu vực của người dùng hay trong một khoảng thời gian. Với cách \textbf{lọc sau}, hệ thống lấy $k$ vector gần nhất trên toàn bộ dữ liệu rồi mới loại những kết quả không thỏa điều kiện, nên có thể trả về ít hơn $k$ kết quả, thậm chí không có kết quả nào, khi khu vực của người dùng chỉ chiếm một phần nhỏ dữ liệu. Với cách \textbf{lọc trước}, hệ thống xác định tập vector thỏa điều kiện rồi chỉ tìm trong tập đó, nên mọi kết quả đều thuộc phạm vi được phép và chỉ ít hơn $k$ khi phạm vi thực sự có ít dữ liệu. Milvus~\cite{wang2021milvus}, hệ quản trị cơ sở dữ liệu vector mã nguồn mở, lưu các trường vô hướng bên cạnh vector và thực hiện tìm kiếm theo cách lọc trước~\cite{milvus2026filtered}.

\subsection{Tìm kiếm vector trong hệ thống}
\label{sec:3.6.4}

Hệ thống dùng Milvus với chỉ mục HNSW, độ đo tích vô hướng, $M = 16$, \textit{ef\-Con\-struc\-tion}~$= 128$ và $\mathit{ef} = 64$. Một lượt tìm kiếm gửi cho Milvus vector truy vấn cùng điều kiện lọc gồm khu vực của người dùng, tập camera được phép, khoảng thời gian nếu có và trạng thái sẵn sàng, rồi nhận về tối đa $k$ kết quả với $k \in \{4, 8, 12, 16\}$. Độ tương đồng được hiển thị dưới dạng điểm phù hợp và chỉ được tính khi truy vấn. Mỗi kết quả được đối chiếu với PostgreSQL để lấy thông tin chi tiết và kiểm tra lại phạm vi quyền; kết quả không còn hợp lệ do hai kho tạm thời chưa đồng bộ bị loại, và nếu còn ít hơn $k$ kết quả, hệ thống tìm lại với số lượng gấp đôi, tối đa ba lượt, nên số lượng cần lấy không vượt quá 64, bằng giá trị $\mathit{ef}$. Với lượng dữ liệu trình diễn của đề tài (1.523 track, Mục~\ref{sec:7.2}), tìm kiếm vét cạn vẫn đủ nhanh; chỉ mục HNSW được chọn để giữ thời gian phản hồi khi dữ liệu tích lũy từ nhiều camera trong thời gian dài.
